%% 02_arquitectura.tex — Motor asíncrono, pipeline de ablación, telemetría.

\section{Arquitectura y diseño}
\label{sec:arquitectura}

\herramienta se estructura en cuatro subsistemas: el núcleo HTTP asíncrono, el
cargador de \emph{payloads}, la jerarquía de \emph{testers} de vulnerabilidad
con su \emph{pipeline} de planificación, y la instrumentación de telemetría y
reproducibilidad. La \cref{fig:arquitectura} resume el flujo de trabajo.

% ---------------------------------------------------------------- figura ----
\begin{figure*}[t]
  \centering
  \resizebox{\textwidth}{!}{%
  \begin{tikzpicture}[node distance=6mm and 10mm]
    % --- entrada
    \node[bloque] (cli) {CLI / \\ lista de objetivos};
    \node[nucleo, right=of cli] (core) {Núcleo HTTP \\ \code{AsyncScannerCore}\\[-1pt]
      \scriptsize sesión \code{aiohttp}, caché GET,\\[-2pt]
      \scriptsize \emph{token bucket}, guarda SSRF};
    \node[nucleo, right=of core] (pool) {\emph{Worker pool} \\ con contrapresión\\[-1pt]
      \scriptsize pares $(tester, url)$};

    % --- pipeline de ablación
    \node[ablac, below=14mm of core] (load) {Cargador de \emph{payloads}\\[-2pt]
      \scriptsize 14 categorías, corpus en disco\\[-2pt]
      \scriptsize \code{PAYLOAD/data/*.json}};
    \node[ablac, right=of load] (sched) {\emph{Pipeline} de planificación\\[-2pt]
      \scriptsize \code{drop\_oob} $\to$ \textbf{entrelazado} $\to$\\[-2pt]
      \scriptsize \textbf{priorización} $\to$ \textbf{orden} $\to$ mutación};
    \node[ablac, right=of sched] (probe) {Sonda + \\ confirmación\\[-2pt]
      \scriptsize reflexión / error / \\[-2pt]
      \scriptsize \code{confirm\_time\_based} ($\tau=\SI{4}{\second}$)};

    % --- artefactos
    \node[artefacto, below=14mm of load] (tel) {\code{TelemetryWorker}\\[-2pt]
      \scriptsize cola no bloqueante,\\[-2pt]
      \scriptsize \code{asyncio.to\_thread} escritura};
    \node[artefacto, right=of tel] (jsonl) {Telemetría JSONL\\[-2pt]
      \scriptsize 1 fila / sonda de \emph{payload}};
    \node[artefacto, right=of jsonl] (manif) {Manifiesto\\[-2pt]
      \scriptsize SHA-256 corpus, \code{git},\\[-2pt]
      \scriptsize semilla global};
    \node[artefacto, right=of manif] (report) {Informe\\[-2pt]
      \scriptsize JSON / HTML,\\[-2pt]
      \scriptsize secretos enmascarados};

    % --- flechas
    \draw[flecha] (cli) -- (core);
    \draw[flecha] (core) -- (pool);
    \draw[flecha] (pool.south) to[out=-90,in=90] (probe.north);
    \draw[flecha] (load) -- (sched);
    \draw[flecha] (sched) -- (probe);
    \draw[flecha] (core.south) to[out=-90,in=90] (load.north);
    \draw[flecha] (probe.south) to[out=-90,in=90] (jsonl.north);
    \draw[flecha] (load.south) to[out=-90,in=90] (tel.north);
    \draw[flecha] (tel) -- (jsonl);
    \draw[flecha] (jsonl) -- (manif);
    \draw[flecha] (manif) -- (report);

    \begin{scope}[on background layer]
      \node[fit=(load)(sched)(probe), draw=orange!50, dashed,
            rounded corners, inner sep=4mm,
            label={[orange!70!black,font=\footnotesize]above:%
            \emph{Pipeline} de ablación de \emph{payloads}}] {};
    \end{scope}
  \end{tikzpicture}%
  }
  \caption{Flujo de trabajo de \herramienta. El núcleo HTTP asíncrono
  concentra toda la E/S de red; el \emph{worker pool} despacha pares
  $(tester, url)$ con contrapresión. Cada \emph{tester} carga su corpus de
  \emph{payloads} y lo hace pasar por el \emph{pipeline} de planificación
  antes de sondear el objetivo. La etapa de confirmación reetiqueta la
  confianza del hallazgo. El \code{TelemetryWorker} escribe una fila JSONL por
  sonda sin bloquear el bucle de eventos; el manifiesto fija los insumos de
  reproducibilidad.}
  \label{fig:arquitectura}
\end{figure*}

\subsection{Núcleo HTTP asíncrono}
\label{sec:nucleo}

Todo el tráfico de red pasa por \code{AsyncScannerCore.request()}, un único
punto de estrangulamiento construido sobre una \code{aiohttp.ClientSession}.
Centralizar la E/S permite aplicar de forma uniforme: (a)~una caché de
respuestas \code{GET} con código \num{200}; (b)~limitación de tasa mediante un
\emph{token bucket}; (c)~rotación de \emph{User-Agent} sembrable; (d)~tiempos
de espera separados para el establecimiento de conexión y la lectura de
\emph{socket}; (e)~un tope de \SI{5}{\mebi\byte} sobre el cuerpo de la
respuesta, con lectura fragmentada y marca de truncamiento; y (f)~una
salvaguarda frente a \emph{Server-Side Request Forgery} que rechaza
redirecciones hacia direcciones privadas salvo autorización explícita. El
campo \code{elapsed} que devuelve el núcleo es monótono y abarca las
redirecciones; para la latencia por sonda con fines científicos se mide en
cambio \code{time.perf\_counter()} en la frontera del \emph{tester}.

El despacho de trabajo emplea \code{run\_worker\_pool()}, que impone
contrapresión mediante una cola acotada (\code{queue\_factor}) de modo que el
descubrimiento de nuevos objetivos no pueda desbordar la memoria ni saturar el
objetivo. La interrupción por \code{SIGINT} es ordenada: se drena la cola y se
cierra la sesión de forma idempotente.

\subsection{\emph{Pipeline} de planificación de \emph{payloads}}
\label{sec:pipeline}

Cada \emph{tester} deriva de una clase base común que resuelve la carga de
\emph{payloads} mediante \code{load\_payloads()}. Internamente esta rutina
aplica, en orden, las siguientes etapas, todas ellas desacoplables por
configuración para el estudio de ablación:

\begin{enumerate}
  \item \code{\_drop\_useless\_oob}: descarta \emph{payloads} de tipo
  \emph{out-of-band} cuando no hay un colector configurado.
  \item \code{\_interleave\_by\_context} (\textsc{entrelazado}): reordena la
  lista para alternar entre contextos sintácticos (HTML, atributo, JavaScript,
  URL, SQL, \dots) en lugar de agotar un contexto antes de pasar al siguiente.
  La intuición es que, bajo un presupuesto reducido, cubrir primero un
  \emph{payload} representativo de cada contexto maximiza la probabilidad de
  acierto temprano.
  \item \code{\_prioritize} (\textsc{priorización}): ordena de forma estable
  por una clave $(\text{confianza \emph{a priori}}, \text{severidad})$, de
  modo que los vectores con mayor probabilidad histórica de éxito y mayor
  impacto se prueben antes.
  \item \code{\_apply\_payload\_order} (\textsc{orden}): política terminal
  configurable --- \code{natural} (identidad), \code{random} (barajado
  sembrado) o \code{freq} (grupos de contexto más poblados primero).
  \item \code{\_apply\_runtime\_mutations}: transformaciones de evasión de WAF
  y variación de caja aplicadas en el momento del envío.
\end{enumerate}

\subsection{Confirmación en tiempo de ejecución y confianza calibrada}
\label{sec:confirmacion}

Cada hallazgo lleva un nivel de confianza de un dominio ordenado
$\{\textsc{low}, \textsc{medium}, \textsc{high}, \textsc{confirmed}\}$. El
corpus asigna una confianza \emph{a priori} a cada \emph{payload}; la etapa de
confirmación produce la confianza \emph{final}:

\begin{itemize}
  \item \textbf{Inyección SQL basada en tiempo.} Se mide la latencia base con
  remuestreo (\code{measure\_baseline\_latency}) y se contrasta la respuesta
  al \emph{payload} temporizado frente a un umbral $\tau = \SI{4}{\second}$
  (\code{confirm\_time\_based}); solo si la diferencia supera el umbral el
  hallazgo se marca \textsc{confirmed}.
  \item \textbf{XSS.} Una reflexión cruda se refina según el contexto HTML en
  el que aparece (\code{\_assess}), promoviendo a \textsc{high} o degradando a
  \textsc{low}.
  \item Bajo la condición de ablación \code{--no-runtime-confirm}, ambas ramas
  se cortocircuitan y el hallazgo se reporta con la confianza \emph{a priori}
  del \emph{payload}, dejando constancia (\code{verification=skipped}) en la
  evidencia. La sonda se registra en telemetría en cualquier caso.
\end{itemize}

\subsection{Telemetría no bloqueante}
\label{sec:telemetria}

El \code{TelemetryWorker} desacopla la generación de eventos de su
persistencia. El método \code{record()} es síncrono y no bloqueante: encola la
fila y asigna \code{request\_index} (correlativo por ejecución), marca de
tiempo ISO-8601 UTC y \rid. Un consumidor asíncrono agrupa las filas en lotes
y las escribe mediante \code{asyncio.to\_thread()}, de forma que el bucle de
eventos nunca se bloquea en disco. Si la cola se llena, las filas excedentes se
descartan y se contabilizan (\code{dropped}), preservando la propiedad de
contrapresión.

La \textbf{granularidad} es deliberada: una fila JSONL por \emph{sonda de
payload}, tras las heurísticas de confirmación del \emph{tester}, y no por
escritura de \emph{socket}. Las subpeticiones de línea base y de confirmación
no se registran individualmente. El esquema de cada fila incluye \code{run\_id,
timestamp, tester\_id, payload\_id, context, confidence\_apriori, url, method,
elapsed\_time, decision, confidence\_final} y \code{request\_index}. Esta traza
es el insumo primario del análisis de coste de detección de la
\cref{sec:resultados}.

\subsection{Reproducibilidad: manifiesto por ejecución}
\label{sec:manifiesto}

Junto a la telemetría, cada ejecución instrumentada emite un manifiesto JSON
que fija: la versión del esquema, el \rid, la marca temporal de creación, el
\emph{commit} de \code{git} (\code{unknown} si no procede), la semilla global
del RNG, y un bloque \code{corpus} con el \emph{hash} SHA-256 calculado sobre
los ficheros \code{PAYLOAD/data/*.json} ordenados (el nombre de cada fichero se
mezcla en el digesto) más la lista de ficheros y su número. Se incluye, además,
el diccionario de configuración completo de la ejecución. El cálculo del
\emph{hash} y la consulta a \code{git} se realizan en un hilo aparte; un fallo
se registra y se ignora, sin abortar ni demorar el escaneo. Este artefacto
permite verificar \emph{a posteriori} que dos ejecuciones comparten exactamente
el mismo corpus y la misma configuración.
